\subsection{实验三：\texttt{TCP} 连接重置脚本}

实验三的脚本监听镜像到 Attacker 的受害者服务器响应报文，并依据其中的 TCP 确认号构造面向受害者服务器的 RST 报文。完整实现见代码清单~\ref{lst:append-task3-rst}。

\lstinputlisting[
    language=Python,
    caption={实验三 TCP RST 注入脚本},
    label={lst:append-task3-rst}
]{../../Experiment3/task3/tcp_reset.py}

\subsection{实验四：\texttt{TCP} 会话劫持脚本}

实验四的脚本等待镜像 Telnet 流量短暂空闲后，使用最近捕获报文中的序列号、确认号和负载长度计算注入报文参数，再发送会话数据。完整实现见代码清单~\ref{lst:append-task4-hijack}。

\lstinputlisting[
    language=Python,
    basicstyle=\ttfamily\footnotesize,
    caption={实验四 Telnet 会话数据注入脚本},
    label={lst:append-task4-hijack}
]{../../Experiment3/task4/topology-lab/tcp_session_hijack.py}
